\chapter{Trois lecteurs : image, son, vidéo}

Un lecteur est l'application la plus simple qui soit utile : on ouvre un fichier,
on l'affiche, on n'écrit rien. C'est aussi celle qui enseigne le plus, parce que
tout le travail est dans la \emph{transformation} --- des pixels, des
échantillons, des images --- et pas dans la logique.

Le dépôt en contient trois. Ils sont réels, ils tournent, et ce chapitre les cite.

\begin{center}
\begin{tabular}{@{}llll@{}}
\toprule
& \textbf{NkImageDemo} & \textbf{NkAudioPlayer} & \textbf{NkVideoPlayer} \\
\midrule
Lignes & 2\,369 & 263 & 684 \\
Modules & NKImage & NKAudio & NKMedia + NKAudio \\
\bottomrule
\end{tabular}
\end{center}

\begin{nkpiege}[À lire avant de copier ces trois-là]
\textbf{Les trois sont antérieurs à \texttt{NkCanvasApp}.} Ils ouvrent la fenêtre
eux-mêmes, dans \texttt{nkmain}, et tiennent leur propre boucle. Le patron
d'aujourd'hui, celui du chapitre sur la coquille, fait tout cela pour vous.

Ce n'est pas une raison de les ignorer --- ils montrent exactement ce que la
coquille vous épargne --- mais c'en est une de ne pas les recopier tels quels.
\textbf{Les convertir sur \texttt{NkCanvasApp} est d'ailleurs le meilleur
exercice de ce chapitre}, et le premier que je vous propose.
\end{nkpiege}

\section{Le lecteur audio : réduire un million d'échantillons à mille colonnes}

Un fichier de trois minutes en 44\,100 Hz contient environ huit millions
d'échantillons. La fenêtre en fait mille de large. Il faut donc en montrer huit
mille par colonne --- et le \emph{comment} n'est pas un détail.

\begin{nkcode}{Applications/NkAudioPlayer/src/main.cpp}
float32 lo = 1.0f, hi = -1.0f;
for (uint64 f = f0; f < f1; ++f) {
    float32 m = 0.0f;
    for (int32 k = 0; k < ch; ++k)
        m += s.data[f * (uint64)ch + (uint64)k];
    m /= (float32)ch;
    if (m < lo) lo = m;
    if (m > hi) hi = m;
}
out[(uint64)c].lo = lo;
out[(uint64)c].hi = hi;
\end{nkcode}

Chaque colonne garde le \textbf{minimum} et le \textbf{maximum} de sa tranche,
puis se dessine comme un trait vertical de l'un à l'autre.

\begin{nkretenir}
\textbf{Une onde sonore oscille autour de zéro.} Si l'on affichait la
\emph{moyenne} de chaque tranche, on obtiendrait une ligne plate au milieu de
l'écran --- techniquement juste, et parfaitement inutile : la moyenne d'un son
fort et la moyenne d'un silence sont toutes deux proches de zéro.

Ce qui porte l'information, c'est l'\textbf{amplitude}, donc l'écart entre le
minimum et le maximum. C'est un cas d'école du piège général : \emph{un agrégat
peut être exact et ne rien décrire.}
\end{nkretenir}

\begin{nkpiege}
Remarquez l'initialisation : \texttt{lo = 1.0f} et \texttt{hi = -1.0f}, à
l'\emph{envers} des bornes attendues. C'est la seule façon correcte d'initialiser
une recherche de minimum et de maximum : la première valeur rencontrée les
remplace tous les deux.

Initialiser à zéro serait faux --- un signal entièrement négatif rendrait
\texttt{hi = 0}, un maximum qui n'existe pas dans les données.

Et la garde qui suit compte autant :
\texttt{if (hi < lo) \{ lo = 0; hi = 0; \}}. Elle traite la \emph{tranche vide} :
sans elle, une colonne sans échantillon garderait \texttt{1} et \texttt{-1}, et
dessinerait un trait pleine hauteur à partir de rien.
\end{nkpiege}

\section{Un rectangle plein, c'est deux triangles}

\begin{nkcode}{Applications/NkAudioPlayer/src/main.cpp}
static void PushQuad(NkVector<NkVertex> &v, float32 x0, float32 y0,
                     float32 x1, float32 y1, const NkColor2D &col) {
    NkVertex a{x0, y0, 0, 0, col.r, col.g, col.b, col.a};
    NkVertex b{x1, y0, 0, 0, col.r, col.g, col.b, col.a};
    NkVertex c{x1, y1, 0, 0, col.r, col.g, col.b, col.a};
    NkVertex d{x0, y1, 0, 0, col.r, col.g, col.b, col.a};
    v.PushBack(a); v.PushBack(b); v.PushBack(c);
    v.PushBack(a); v.PushBack(c); v.PushBack(d);
}
\end{nkcode}

\begin{nkretenir}
\textbf{Quatre coins, six sommets.} Le GPU ne connaît que des triangles : un
rectangle est le triangle $abc$ plus le triangle $acd$. Les sommets $a$ et $c$
sont donc envoyés \emph{deux fois} --- c'est normal, c'est le prix d'une liste de
triangles.

C'est exactement ce que fait \texttt{AddRectFilled} pour vous. Le voir écrit une
fois explique tout le reste : pourquoi un dessin coûte en \emph{nombre de
sommets}, et pourquoi regrouper les rectangles dans un même tampon est ce qui
rend un rendu rapide.
\end{nkretenir}

\section{Le lecteur vidéo : la chaîne complète}

Son en-tête donne le patron en six lignes, et il vaut d'être cité tel quel :

\begin{nkcode}{Applications/NkVideoPlayer/src/main.cpp}
//   1. NkWindow          -> fenetre native (Win/Linux/macOS/...)
//   2. NkRenderWindow    -> cible de rendu 2D NKCanvas
//   3. NkVideoReader     -> decode chaque image en RGBA
//   4. NkTexture         -> recoit les pixels RGBA de l'image courante (Update)
//   5. NkSprite          -> affiche la texture, mise a l'echelle
//   6. boucle cadencee au fps de la video (Clear -> Draw -> Display)
\end{nkcode}

\begin{nkretenir}
\textbf{Il n'y a aucune API « vidéo » dans le rendu.} Une vidéo, pour l'écran,
c'est une texture qu'on remplace à chaque image. \texttt{frameTex.Update(...)}
téléverse les nouveaux pixels, le sprite les affiche, et la boucle recommence.

Le décodage --- H.264, MJPEG, conteneurs --- est entièrement dans NKMedia, et le
rendu n'en sait rien. C'est la séparation la plus profitable de tout ce cours :
\emph{le module qui produit des pixels ignore le module qui les montre}.
\end{nkretenir}

\begin{nkcode}{Applications/NkVideoPlayer/src/main.cpp}
frameTex.Update(fr.rgba.Data(), (uint32)fr.width, (uint32)fr.height, 0, 0);
\end{nkcode}

Une seule ligne, appelée une fois par image. Tout le reste du fichier est de la
cadence, de la navigation et de l'étalonnage de couleur.

\section{Traiter les pixels avant de les téléverser}

Le lecteur applique un étalonnage --- noir et blanc, sépia, table de correspondance
\texttt{.cube} --- et l'endroit où il le fait est un enseignement à lui seul :

\begin{nkcode}{Applications/NkVideoPlayer/src/main.cpp (commentaire d'origine)}
// cablage est le meme : transformer les pixels juste avant frameTex.Update
\end{nkcode}

\begin{nkretenir}
\textbf{Un traitement d'image se branche entre le décodeur et la texture.} Ni
avant (le décodeur ne doit rien savoir de vos couleurs), ni après (une fois
téléversés, les pixels sont sur le GPU et ne vous appartiennent plus).

Cette position est un \emph{point d'extension} : ajouter un filtre ne demande de
toucher à rien d'autre.
\end{nkretenir}

\begin{nkpiege}[La dette qui est dans ce fichier, et qu'il faut voir]
\texttt{NkVideoPlayer} inclut encore \texttt{<cstdio>}, \texttt{<cstring>} et
\texttt{<cstdlib>}.

C'est contraire à la règle du dépôt --- \emph{aucun \texttt{std::}, pas même les
fonctions C} --- et ce n'est pas un oubli anodin : les équivalents maison
existent tous (\texttt{NkFile}, \texttt{NkStringUtils}, \texttt{NkFormat}). Le
fichier est plus ancien que la règle.

Je le montre plutôt que de le cacher, parce que c'est ce que vous rencontrerez :
\textbf{un dépôt réel contient toujours du code antérieur à ses propres règles}.
Ce qui compte n'est pas qu'il n'y en ait jamais, c'est que ce soit \emph{nommé}.
\end{nkpiege}

\section{Le lecteur d'images, et ce qu'il ne faut pas en retenir}

\texttt{NkImageDemo} fait 2\,369 lignes pour afficher une image --- neuf fois le
lecteur audio. La raison est visible dans ses fichiers :

\begin{center}
\begin{tabular}{@{}ll@{}}
\toprule
\textbf{Fichier} & \textbf{Ce que c'est} \\
\midrule
\texttt{Render/GLContext} & un contexte OpenGL, à la main \\
\texttt{Render/GLRenderer2D} & un renderer 2D, à la main \\
\texttt{Render/Texture2D} & une texture, à la main \\
\texttt{Render/FontAtlas} & un atlas de police, à la main \\
\texttt{Render/SafeArea} & la zone sûre, à la main \\
\texttt{ViewerApp} & le visualiseur proprement dit \\
\bottomrule
\end{tabular}
\end{center}

\begin{nkretenir}
\textbf{Cinq de ses six fichiers réimplémentent ce que NKCanvas porte déjà.} Le
visualiseur lui-même --- charger, zoomer, déplacer, afficher --- est la plus
petite partie.

C'est le doublon dans sa forme la plus coûteuse : le code marche, personne n'a
tort, et pourtant chaque correction apportée à \texttt{NkRenderWindow} ne
parviendra jamais à \texttt{GLRenderer2D}.

\textbf{La leçon} : avant d'écrire un mécanisme, chercher qui le porte déjà --- et
commencer par la couche du dessous. Un \texttt{grep} lancé dans son propre
dossier ne trouve jamais ce que la bibliothèque porte.
\end{nkretenir}

\begin{nkbilan}
Un lecteur ne décide de rien : tout son travail est une transformation. Réduire
un signal pour l'afficher demande le \emph{minimum et le maximum}, jamais la
moyenne --- une moyenne exacte peut ne rien décrire. Un rectangle plein est deux
triangles et six sommets. Une vidéo, pour le rendu, n'est qu'une texture qu'on
remplace à chaque image, et le traitement des pixels se branche entre le décodeur
et la texture. Enfin, \texttt{NkImageDemo} montre ce que coûte un doublon :
2\,369 lignes dont cinq sixièmes existaient déjà un étage plus bas.
\end{nkbilan}

\begin{nkquestions}
\begin{enumerate}
\item Pourquoi une forme d'onde s'affiche-t-elle en min/max et non en moyenne ?
\item Pourquoi \texttt{lo} est-il initialisé à $1$ et \texttt{hi} à $-1$ ?
\item Combien de sommets pour un rectangle plein, et pourquoi pas quatre ?
\item Où se branche un traitement d'image dans la chaîne du lecteur vidéo, et
      pourquoi pas ailleurs ?
\item Pourquoi \texttt{NkImageDemo} est-il neuf fois plus gros que
      \texttt{NkAudioPlayer} ?
\end{enumerate}
\end{nkquestions}

\begin{nkpratique}
Réécrivez \texttt{NkAudioPlayer} sur \texttt{NkCanvasApp}.

Vous devez obtenir le même résultat à l'écran --- la forme d'onde, le curseur de
lecture --- avec \textbf{aucun appel à \texttt{NkWindow} ni à
\texttt{NkRenderWindow}} dans votre code. Comptez les lignes avant et après, et
dites lesquelles ont disparu et pourquoi.

Puis faites de même pour le lecteur vidéo.
\end{nkpratique}

\begin{nkdemo}
Prouvez que la moyenne ne convient pas.

Ajoutez à votre lecteur un second mode d'affichage qui dessine la \emph{moyenne}
de chaque tranche au lieu du min/max, et une touche pour basculer entre les deux.
Chargez un morceau réel et capturez les deux images.

\textbf{Ce que la démonstration doit établir} : non pas que la seconde est
« moins jolie », mais qu'elle est \emph{plate} --- et donc qu'elle ne distingue
pas un passage fort d'un passage faible. Mesurez-le : affichez l'écart entre la
plus grande et la plus petite valeur affichée, dans les deux modes.
\end{nkdemo}

\begin{nklogique}
Une vidéo à 25 images par seconde dure 4 minutes. Le décodage d'une image prend
en moyenne 12 ms, le téléversement 2 ms, le dessin 1 ms.

\begin{enumerate}
\item La lecture tiendra-t-elle la cadence ? Justifiez par le calcul.
\item Une seule image sur cinquante prend 300 ms à décoder (image-clé).
      La lecture saccade-t-elle ? Que faudrait-il changer, sans accélérer le
      décodeur ?
\item Vous ajoutez un filtre couleur qui coûte 9 ms par image. Où le mettre ---
      et si vous n'aviez le droit de le faire que sur une image sur deux, la
      vidéo serait-elle acceptable ? Argumentez sur ce que voit l'\oe il.
\end{enumerate}
\end{nklogique}
